% ============================================================================
% LEHRERHINWEISE: Die Gewichtskraft (UE 3)
% Physik Klasse 7 – Gesamtschule MV – 45 Minuten
% ============================================================================
\documentclass[a4paper,11pt]{article}

\usepackage[left=2cm,right=2cm,top=2cm,bottom=2cm]{geometry}
\usepackage[utf8]{inputenc}
\usepackage[T1]{fontenc}
\usepackage[ngerman]{babel}
\usepackage{fancyhdr}
\usepackage{tcolorbox}
\usepackage{amssymb,amsmath}
\usepackage{xcolor}
\usepackage{array}
\usepackage{tabularx}
\usepackage{booktabs}
\usepackage{enumitem}
\usepackage{siunitx}
\sisetup{locale=DE, per-mode=symbol}

\renewcommand{\familydefault}{\sfdefault}

\pagestyle{fancy}
\fancyhf{}
\renewcommand{\headrulewidth}{0.4pt}
\lhead{\textbf{Lehrerhinweise}}
\chead{Die Gewichtskraft (UE 3)}
\rhead{Physik Kl. 7}
\setlength{\headheight}{14pt}
\cfoot{\thepage}

\definecolor{physik}{HTML}{2c5aa0}
\definecolor{hinweis}{HTML}{b35c00}
\definecolor{fehler}{HTML}{c62828}

\newcommand{\stern}{$\bigstar$}

\setlength{\parindent}{0pt}

\begin{document}

% ==================== ÜBERSICHT ====================

\section*{Stundenziel}

Die Schülerinnen und Schüler berechnen die Gewichtskraft mit der Formel $F_G = m \cdot g$ und erklären die Ortsabhängigkeit der Gewichtskraft.

\vspace{3mm}

\textbf{Feinziele:}
\begin{itemize}[nosep]
    \item FZ1: SuS nennen die Definition der Masse und der Gewichtskraft (AFB I)
    \item FZ2: SuS berechnen $F_G$ mit der Formel $F_G = m \cdot g$ (AFB II)
    \item FZ3: SuS wenden die Umrechnung \SI{100}{\gram} $\hat{=}$ \SI{1}{\newton} an (AFB I/II)
    \item FZ4: SuS erklären die Ortsabhängigkeit von $g$ (AFB II)
    \item FZ5: SuS unterscheiden Masse und Gewichtskraft (AFB II/III)
\end{itemize}

\vspace{5mm}

% ==================== STUNDENVERLAUF ====================

\section*{Stundenverlauf (45 Min)}

\begin{tabularx}{\textwidth}{|c|l|X|c|c|}
\hline
\textbf{Min} & \textbf{Phase} & \textbf{Inhalt} & \textbf{SF} & \textbf{Material} \\
\hline
0--5 & Einstieg & Frage: ,,Warum fallen Dinge runter?'' $\rightarrow$ Erde zieht an & PL & -- \\
\hline
5--10 & Erarbeitung I & Masse einführen: Merksatz 1 ausfüllen (Stoff, kg, überall gleich) & EA & LP 1--2, AB \\
\hline
10--18 & Erarbeitung II & Gewichtskraft einführen: $F_G = m \cdot g$, $g = \SI{10}{\newton\per\kilogram}$ & EA & LP 3--4, AB \\
\hline
18--22 & Übung I & Einfache Berechnungen: $F_G$ berechnen (Aufg. 1) & EA & LP 5, AB \\
\hline
22--28 & Übung II & Faustformel \SI{100}{\gram} $\hat{=}$ \SI{1}{\newton}, Umrechnung (Aufg. 2+3) & EA & LP 6, AB \\
\hline
28--33 & Vertiefung & Ortsabhängigkeit: Mond, Mars, Jupiter (Aufg. 4+5) & EA & LP 7, AB \\
\hline
33--45 & Stundenleistung & Digitaler Minitest ($\sim$26 BE) & EA & MT \\
\hline
\end{tabularx}

\vspace{5mm}

% ==================== DIFFERENZIERUNG ====================

\section*{Differenzierung}

\begin{tabularx}{\textwidth}{|c|X|}
\hline
\textbf{Niveau} & \textbf{Aufgabenbereich} \\
\hline
\stern{} Basis & Aufgabe 1 + 2: Formel anwenden mit ganzen Zahlen, Faustformel \\
\hline
\stern\stern{} Standard & Aufgabe 3 + 4: Dezimalzahlen, Umrechnung g $\leftrightarrow$ kg, Formel umstellen nach $m$ \\
\hline
\stern\stern\stern{} Erweitert & Aufgabe 5 + 6: Ortsabhängigkeit berechnen, Masse vs. Gewicht erklären \\
\hline
\end{tabularx}

\vspace{3mm}

\textbf{Minitest:} Schüler wählen selbst ihr Niveau (\stern/\stern\stern/\stern\stern\stern). Die Bewertung erfolgt differenziert:
\begin{itemize}[nosep]
    \item \stern: max. 16 BE erreichbar
    \item \stern\stern: max. 22 BE erreichbar
    \item \stern\stern\stern: max. 26 BE erreichbar
\end{itemize}

\vspace{5mm}

% ==================== WICHTIGER HINWEIS ====================

\section*{Wichtiger Hinweis: Der Ortsfaktor $g$}

\begin{tcolorbox}[colback=hinweis!10,colframe=hinweis,title=Ortsfaktor / Fallbeschleunigung / Gravitationsfeldstärke]
\textbf{Exakter Wert:} $g = \SI{9,81}{\newton\per\kilogram}$ (auf der Erde)

\textbf{Wir rechnen mit:} $g = \SI{10}{\newton\per\kilogram}$ (Abweichung nur ca. 2\%)

\vspace{2mm}
\textbf{Begründung für die Vereinfachung:}
\begin{itemize}[nosep]
    \item Kopfrechenbar $\rightarrow$ mehr Fokus auf das Verstehen der Physik
    \item Keine Ablenkung durch Taschenrechner
    \item Bei realitätsnahen Anwendungen später: exakter Wert verwenden
\end{itemize}

\vspace{2mm}
\textbf{Alternative Bezeichnungen:} Der Ortsfaktor $g$ wird auch \textbf{Fallbeschleunigung} oder \textbf{Gravitationsfeldstärke} genannt.
\end{tcolorbox}

\vspace{5mm}

% ==================== HÄUFIGE FEHLER ====================

\section*{Häufige Fehler und Reaktionen}

\begin{tabularx}{\textwidth}{|X|X|}
\hline
\textbf{Fehler} & \textbf{Reaktion / Korrektur} \\
\hline
$g = 9{,}81$ statt 10 & ,,Wir rechnen mit $g = \SI{10}{\newton\per\kilogram}$ -- kopfrechenbar!'' (Abweichung nur 2\%) \\
\hline
Masse und Gewicht verwechselt & ,,Masse ist die Stoffmenge (\si{\kilogram}), Gewichtskraft ist die Anziehungskraft (\si{\newton}). Merksatz 1 nochmal lesen!'' \\
\hline
Einheit vergessen & ,,Ergebnis immer mit Einheit! $F_G$ hat die Einheit Newton (\si{\newton})'' \\
\hline
Bei \SI{500}{\gram}: $F_G = 500 \cdot 10 = \SI{5000}{\newton}$ & ,,Erst in \si{\kilogram} umrechnen! \SI{500}{\gram} = \SI{0,5}{\kilogram} $\Rightarrow F_G = \SI{5}{\newton}$'' \\
\hline
Formel nicht umstellen können & Formeldreieck zeigen: $F_G$ oben, $m$ und $g$ unten \\
\hline
,,Auf dem Mond wiegt man weniger'' & ,,Korrekter: Die Gewichtskraft ist kleiner. Die Masse bleibt gleich!'' \\
\hline
\end{tabularx}

\vspace{5mm}

% ==================== MATERIAL ====================

\section*{Benötigtes Material}

\begin{itemize}[nosep]
    \item Tablets/Laptops für Lernpfad und Minitest
    \item Arbeitsblatt AB-03-gewichtskraft (1× pro Schüler)
    \item Optional: Federkraftmesser + verschiedene Massen zur Demonstration
    \item Optional: Tafelbilder Formel + Formeldreieck
\end{itemize}

\vspace{5mm}

% ==================== VORWISSEN ====================

\section*{Vorwissen (aus UE 1 + 2)}

\begin{itemize}[nosep]
    \item Kraftwirkungen kennen (Bewegungsänderung, Verformung)
    \item Kraft als physikalische Größe: $F$, Newton (N), Federkraftmesser
    \item Kraftpfeil: Angriffspunkt, Richtung, Betrag
    \item Kräfte zusammensetzen ($F_{\text{res}}$)
\end{itemize}

\vspace{5mm}

% ==================== LÖSUNGEN ====================

\section*{Lösungen zum Arbeitsblatt}

\textbf{Merksatz 1 (Masse):}
\begin{itemize}[nosep]
    \item ,,\ldots wie viel \textbf{Stoff} in einem Körper enthalten ist.''
    \item Einheit: \textbf{kg} (Kilogramm)
    \item Die Masse ist \textbf{überall gleich / konstant}
\end{itemize}

\vspace{3mm}

\textbf{Merksatz 2 (Gewichtskraft):}
\begin{itemize}[nosep]
    \item ,,\ldots einen Körper \textbf{anzieht}.''
    \item $F_G$ in \textbf{\si{\newton} (Newton)}, $m$ in \textbf{\si{\kilogram}}, $g$ in \textbf{\si{\newton\per\kilogram}}
\end{itemize}

\vspace{3mm}

\textbf{Aufgabe 1:}
\begin{itemize}[nosep]
    \item a) $F_G = \SI{5}{\kilogram} \cdot \SI{10}{\newton\per\kilogram} = \SI{50}{\newton}$
    \item b) $F_G = \SI{8}{\kilogram} \cdot \SI{10}{\newton\per\kilogram} = \SI{80}{\newton}$
    \item c) $F_G = \SI{20}{\kilogram} \cdot \SI{10}{\newton\per\kilogram} = \SI{200}{\newton}$
\end{itemize}

\vspace{3mm}

\textbf{Aufgabe 2:}
\begin{itemize}[nosep]
    \item a) \SI{200}{\gram} $\rightarrow F_G = \SI{2}{\newton}$
    \item b) \SI{500}{\gram} $\rightarrow F_G = \SI{5}{\newton}$
    \item c) \SI{1000}{\gram} $\rightarrow F_G = \SI{10}{\newton}$
\end{itemize}

\vspace{3mm}

\textbf{Aufgabe 3:}
\begin{itemize}[nosep]
    \item a) $F_G = \SI{2,5}{\kilogram} \cdot \SI{10}{\newton\per\kilogram} = \SI{25}{\newton}$
    \item b) $F_G = \SI{0,4}{\kilogram} \cdot \SI{10}{\newton\per\kilogram} = \SI{4}{\newton}$
    \item c) \SI{750}{\gram} = \textbf{\SI{0,75}{\kilogram}} $\rightarrow F_G = \SI{0,75}{\kilogram} \cdot \SI{10}{\newton\per\kilogram} = \SI{7,5}{\newton}$
\end{itemize}

\vspace{3mm}

\textbf{Aufgabe 4:}
\begin{itemize}[nosep]
    \item a) $m = \dfrac{\SI{40}{\newton}}{\SI{10}{\newton\per\kilogram}} = \SI{4}{\kilogram}$
    \item b) $m = \dfrac{\SI{150}{\newton}}{\SI{10}{\newton\per\kilogram}} = \SI{15}{\kilogram}$
\end{itemize}

\vspace{3mm}

\textbf{Aufgabe 5:} (Astronaut mit $m = \SI{80}{\kilogram}$)
\begin{itemize}[nosep]
    \item Erde: $F_G = \SI{80}{\kilogram} \cdot \SI{10}{\newton\per\kilogram} = \SI{800}{\newton}$
    \item Mond: $F_G = \SI{80}{\kilogram} \cdot \SI{1,6}{\newton\per\kilogram} = \SI{128}{\newton}$
    \item Mars: $F_G = \SI{80}{\kilogram} \cdot \SI{4}{\newton\per\kilogram} = \SI{320}{\newton}$
    \item Jupiter: $F_G = \SI{80}{\kilogram} \cdot \SI{25}{\newton\per\kilogram} = \SI{2000}{\newton}$
\end{itemize}

\vspace{3mm}

\textbf{Aufgabe 6:}
\begin{itemize}[nosep]
    \item a) Die Masse gibt an, wie viel Stoff in einem Körper enthalten ist. Diese Stoffmenge ändert sich nicht, egal wo sich der Körper befindet.
    \item b) Die Gewichtskraft hängt vom Ortsfaktor $g$ ab. Der Ortsfaktor ist auf verschiedenen Himmelskörpern unterschiedlich groß, weil sie unterschiedliche Massen haben.
\end{itemize}

\end{document}
